%####################################################################
% 上智大学情報理工系修士論文・卒業論文テンプレート
% Version 2023.1.30
% Author: Leo Liu @ Yairi Lab.

% 設定箇所は２箇所あります（環境設定・追加設定）．
% citesort は通常の latex に含まれていないので，ファイルと一緒に配布しています.
% 追加設定に，便利なパッケージやマクロを既に追加されています．

% Overleafを使用する場合，Compiler をデフォルトの「pdfLaTeX」でなく，
% 「LaTeX」に設定した上で使ってください．

% 英語の場合，2段組構成はお勧めです．
%####################################################################


%############################## 環境設定 ##############################

% 英語で執筆する場合は \ENGtrue の前の % を削除
\newif\ifENG
% \ENGtrue

% 定理型環境の末尾の□を消したい場合は \QEDtrue の前に % を入れる
\newif\ifQED
% \QEDtrue

% 2段組にする場合は，\TWOCtrue の前の % を削除
\newif\ifTWOC
% \TWOCtrue

% サブタイトルを無効にする場合は，\SUBtrue の前に % を入れる
\newif\ifSUB
\SUBtrue

% サンプルモード．本文の執筆を始める前に \SAMtrue の前に % を入れる
\newif\ifSAM
\SAMtrue

% コンパクトモード．無効にする場合は，\CMPtrue の前に % を入れる
% 要旨を執筆する際はコンパクトモードはおすすめです．
% コンパクトモードを無効にすると，レポートクラスのデフォルトページサイズになります．
\newif\ifCMP
\CMPtrue

%######################### 環境設定, ここまで #########################


%==================== 以下修正不要 =========================
%
% 内容を理解したうえで修正する分にはかまいません。(自己責任で)
%

%**********************************************************
% ドキュメントクラス設定
%
\ifENG
    \ifTWOC
        \documentclass[a4paper,twocolumn]{report}
    \else
        \documentclass[a4paper]{report}
    \fi
\else
    \ifTWOC
        \documentclass[a4paper,twocolumn]{jreport}
    \else
        \documentclass[a4paper]{jreport}
    \fi
    \def\baselinestretch{.9}
\fi
%**********************************************************

%**********************************************************
% ページサイズ設定
%
\ifCMP
    \textheight=252mm
    \textwidth=170mm
    \columnsep=10mm
    
    \voffset=-1in
    \headheight=13mm
    \headsep=12mm
    \topmargin=0mm
    
    \hoffset=-1in
    \oddsidemargin=20mm
    \evensidemargin=20mm
\fi
%**********************************************************

%**********************************************************
% 定理型環境
%
\usepackage{theorem}
\theorembodyfont{\normalfont}

\makeatletter
\def\QED{{\unskip\nobreak\hfil\penalty50
\hskip1em\hbox{}\nobreak\hfil $\Box$\parfillskip\z@
\finalhyphendemerits\z@\par}}
\makeatother

\ifENG
\def\Definition{Definition}
\else
\def\Definition{定義}
\fi

\ifQED
\newtheorem{defi}{\Definition}
\newenvironment{definition}{\begin{defi}}{\QED\end{defi}}
\else
\newtheorem{definition}{\Definition}
\fi

\ifENG
\def\Theorem{Theorem}
\else
\def\Theorem{定理}
\fi

\ifQED
\newtheorem{theo}{\Theorem}
\newenvironment{theorem}{\begin{theo}}{\QED\end{theo}}
\else
\newtheorem{theorem}{\Theorem}
\fi

\ifENG
\def\Proposition{Proposition}
\else
\def\Proposition{命題}
\fi

\ifQED
\newtheorem{prop}{\Proposition}
\newenvironment{proposition}{\begin{prop}}{\QED\end{prop}}
\else
\newtheorem{proposition}{\Proposition}
\fi
%**********************************************************

%==================== 修正不要, ここまで ====================


%############################## 追加設定 ##############################

%**********************************************************
% 使用パッケージ
%
\usepackage[utf8]{inputenc}

% IEEE基準
\usepackage{amsmath,amssymb,graphicx}
\usepackage{algorithmic,amsfonts,cite,graphicx,textcomp,ulem}
\usepackage{xcolor,tabularx}

% サンプルファイル用
\usepackage{dtk-logos}
\usepackage{listings}
\usepackage{textcomp}
%**********************************************************

%**********************************************************
% マクロ定義
%
\newcommand{\etal}{\textit{et al}. }
\newcommand{\ie}{\textit{i}.\textit{e}. }
\newcommand{\eg}{\textit{e}.\textit{g}. }
\newcommand{\etc}{\textit{etc}. }
%**********************************************************

%######################### 追加設定，ここまで #########################


%%%%%%%%%%%%%%%%%%%%%%%%%%%%%%%%%%%%%%%%%%
\begin{document}

%%%%%%%%%%%%%%%%%%%%%%%%%%%%%%%%%%%%%%%%%%
\begin{titlepage}
    \begin{center}
        \Large 2022年度　学位論文 \\

        \vfill
        \Huge\textbf{論文のすゝめ}\\
        \ifSUB
            \huge\textbf{For Theses, With Theses}\\
        \fi

        \vfill
        \Large 理工学研究科理工学専攻情報学領域 \\
        \medskip
        \Large 博士前期課程 \\
        \vspace{\baselineskip}
        \Large \textbf{B1234567} \\
        \smallskip
        \Large \textbf{上智花子} \\
        \vspace{\baselineskip}
        \large 指導教員 \quad 上智太郎\\

        \vspace{3\baselineskip}
        \number\year 年\number\month 月\number\day 日
    \end{center}

    \ifSAM
        \footnotetext[1]{Sample Mode.}
        \footnotetext[2]{Copyright 2022-2023 Leo Liu. All Rights Reserved.}
    \fi
\end{titlepage}

%%%%%%%%%%%%%%%%%%%%%%%%%%%%%%%%%%%%%%%%%%
\cleardoublepage
\pagenumbering{gobble}
\tableofcontents
\cleardoublepage
\pagenumbering{arabic}

%%%%%%%%%%%%%%%%%%%%%%%%%%%%%%%%%%%%%%%%%%
\chapter{はじめに}

このファイルは \LaTeX を利用して情報理工系修士論文・卒業論文を作成するためのサンプルファイルである．本ファイルに記述してある指示に従い，その中身を書き換えることで要旨を作成することができる．

Overleaf で使用する場合，コンパイラをデフォルトの \pdflatex ではなく，\LaTeX に設定した上で使ってください．

%%%%%%%%%%%%%%%%%%%%%%%%%%%%%%%%%%%%%%%%%%
\chapter{論文の体裁}

\section{ページ設定}

基本は横書きで記述し，分量は各教員・研究科・学科の指示に従うこと．２種類のページサーズが用意されている．

\begin{itemize}
    \item \textbf{コンパクトモード}：用紙はA4サイズ(210mm×297mm)縦置き，ページ余白は上：25mm，左右・下：20mmとする．余白は少なめ．要旨などに．
    \item \textbf{デフォルトモード}：\lstinline{report} または \lstinline{jreport} クラスのデフォルトページサイズ．余白は多め．製本用論文正文などに．
\end{itemize}

\section{本文}

2段組みか，1段組みで執筆すること．詳細は各教員・研究科・学科の指示に従うこと．

\section{フォントサイズ}

フォントサイズは 10pt に設定されてあります．

\section{文献の引用}

文献の引用は，\lstinline{thebibliography} 環境または \BibTeX を利用できる．

引用した文献は参考文献に明記し，引用箇所には\cite{F}, \cite{H-K}, \cite{F,H-K,Z}のように文献番号を表記すること．参考文献の表記方法については指導教員の指示に従うこと．

\section{図・表について}

表の挿入は，\lstinline{tabular} または \lstinline{tabularx} 環境を利用できる．

図・表には図\ref{Figure_1}，表\ref{Table_1} のように図表番号をつけ，タイトルとともに，図は下部，表は上部に記すこと.

\begin{table}[h]
    \begin{center}
        \caption{表のタイトル}
        \label{Table_1}
        \medskip
        \begin{tabular}{|l|c|r|} \hline
            \multicolumn{3}{|c|}{\textgt{請求書}} 		\\ \hline
            \multicolumn{1}{|c|}{品名} & 数量 & 金額  	\\ \hline
            \textit{The \TeX{}book}    &    1 & \$30.00 	\\ 
            \textit{\TeX: The Program} &    1 & \$43.00 	\\ \hline
        \end{tabular}
    \end{center}
\end{table}

\begin{figure}[htbp]
    \begin{center}
        \includegraphics[width=40mm,clip]{graph.eps}
        \caption{図のタイトル}
        \label{Figure_1}
    \end{center}
\end{figure}

\section{数式について}

数式には必要に応じて式番号をつけること．数式の例を以下に示す．

\begin{equation}
 H = -\sum_{i=1}^{n} p_i \log p_i 
\end{equation}

\begin{gather}
 E(X) = \sum_{i=1}^{n} x_i P(X=x_i) \\
 V(X) = E((X-E(X))^2) = E(X^2)-(E(X))^2
\end{gather}

%%%%%%%%%%%%%%%%%%%%%%%%%%%%%%%%%%%%%%%%%%
\chapter{本ファイルの利用について}

\section{執筆言語}

論文を英語で記述する場合は, 本ファイル先頭の「環境設定」の中にある指示に従って処理すること．

\section{論文題目・学生氏名}

論文題目・学生氏名は，本文の \lstinline{titlepage} 部分に書き入れること．また，本文を英語で書く場合，題目・学生氏名とも英語で記載すること（ただし，指示に従うこと．）．なお, 副題（英語論文の日本語タイトル，またはその逆）も環境設定で設定の上，\lstinline{titlepage} 部分に書き入れること．

\section{定義・定理などの定理型環境と末尾の処理について}

定義(Definition)，定理(Theorem)，命題(Proposition)，補題(Lemma)，系(Corollary)，注意(Remark)，例(Example)などは定理環境を利用すること．

本ファイルでは 「定義・定理・命題」のマクロを用意してあるので，必要があればそれを利用すること．また，これら3つの定理型環境の末尾に $\Box$ を出力するマクロを用意してあるので，本ファイルの環境設定の指示に従って有効・無効を切り替えること．

独自に作成した定理型環境内で $\Box$ を出力するマクロを使用する方法は，本ファイル先頭にある「修正不要」の中の「定理型環境」を参考にすること．

\begin{definition}
    定理型環境(definition)の利用
\end{definition}

\begin{theorem}
    定理型環境(theorem)の利用
\end{theorem}

\begin{proposition}
    定理型環境(proposition)の利用
\end{proposition}

%%%%%%%%%%%%%%%%%%%%%%%%%%%%%%%%%%%%%%%%%%
\chapter{おわりに}

本稿では，\LaTeX により論文を作成する際の留意事項について解説した．ここに記載されていないことについては指導教員に確認し，その指示に従うこと．

%%%%%%%%%%%%%%%%%%%%%%%%%%%%%%%%%%%%%%%%%%
\cleardoublepage
\pagenumbering{gobble}
\onecolumn

%=====================================
% References, variant A: external bibliography
%=====================================
%\bibliography{your_external_BibTeX_file}

%=====================================
% References, variant B: internal bibliography
%=====================================
\begin{thebibliography}{9}
    % Reference 1
    \bibitem{F} W.~Fulton, ``Young tableaux,'' London Mathematical Society Student Text {\bf 35}, Cambridge, 1996.
    % Reference 2
    \bibitem{H-K} J.~Hong and S-J~Kang, ``Introduction to Quantum Groups and Crystal Bases,'' American Mathematical Society, 2002.
    % Reference 3
    \bibitem{Z} A.~V.~Zelevinsky, ``A Generalization of the Littlewood-Richardson Rule and the Robinson-Shensted-Knuth Correspondence,'' J.~Math. {\bf 69}, pp.82--94, 1981.
\end{thebibliography}

%%%%%%%%%%%%%%%%%%%%%%%%%%%%%%%%%%%%%%%%%%
\end{document}
